\subsection{A finite coefficient calculation, without differentiating a limit}
\label{sec:finite-jet-closure}

The identity $c_j=\varepsilon+h_j$ deserves an explicit arithmetic
justification.  The following calculation replaces coefficient comparison
and repeated differentiation of infinite series by finite convolution and
finite summation by parts.  Throughout this subsection logarithms stand for
the rational names already specified; each displayed identity is verified
to an arbitrarily assigned rational accuracy by a finite calculation.

Put
\[
 b_j(t)=\frac{(\log t)^j}{j!t^2},\qquad
 Z_j(N)=\sum_{n=1}^N b_j(n),\qquad
 M_j(N)=\sum_{n=1}^N\chi(n)b_j(n).
\]
The convention is $(\log1)^0=1$.  Let $Z_j$ and $M_j$ denote the corresponding
names with their stated tails.  The elementary exponential inequality gives
\[
 \frac{(\log t)^j}{j!}\leq4^j t^{1/4}\qquad(t\geq1),
\]
and hence, for both moment sums,
\begin{equation}
 |Z_j-Z_j(N)|,\ |M_j-M_j(N)|
       \leq\frac43\,4^jN^{-3/4}
       \leq2\,4^jN^{-3/4}.
 \label{eq:moment-series-tail}
\end{equation}
These are bounds on ordinary positive tails, so they apply to complex
characters without cancellation.  The coarser estimate
\[
 |Z_j|,\ |M_j|,\ |Z_j(N)|,\ |M_j(N)|\leq3\,2^j
\]
follows from $(\log t)^j/j!\leq2^j\sqrt t$ and
$\sum_{n\geq1}n^{-3/2}\leq3$.

\paragraph{Finite summation by parts identifies the integral coefficients.}
For each integer $N\geq1$ one has the finite identity
\begin{equation}
 M_j(N)=A_\chi(N)b_j(N)-\int_1^N A_\chi(t)b'_j(t)\,dt.
 \label{eq:finite-jet-partial-summation}
\end{equation}
Indeed, split at the integers and telescope
\[
 \sum_{n=1}^{N-1}A_\chi(n)\bigl(b_j(n+1)-b_j(n)\bigr).
\]
This is a bounded double-sum identity.  On each unit interval the formula
for $b'_j$ is checked by the elementary product rule and the finite
logarithm difference estimates.  Replacing its integral by rectangles
leaves the explicitly bounded quadrature error; thus this argument does
not presume an infinitary summation-by-parts theorem.

For $j\geq1$,
\[
 -b'_j(t)=t^{-3}\left(\frac{2(\log t)^j}{j!}
            -\frac{(\log t)^{j-1}}{(j-1)!}\right),
\]
and $-b'_0(t)=2t^{-3}$.  These are exactly the integrands defining $l_j$.
The endpoint term in~\eqref{eq:finite-jet-partial-summation} is bounded by
\[
 |A_\chi(N)b_j(N)|\leq q4^jN^{-7/4}.
\]
Together with the moment tail and the integral tails previously given,
this proves $M_j=l_j$ to any requested accuracy, with primitive recursive
cutoffs.

Apply the same finite identity with the constant sequence $1$ in place of
$\chi$; its partial-sum function is $\lfloor t\rfloor=t-\{t\}$.  One more
finite integration by parts gives
\begin{equation}
 Z_j(N)=\mathbf1_{j=0}+\int_1^N b_j(t)\,dt
                         +\int_1^N\{t\}b'_j(t)\,dt.
 \label{eq:finite-zeta-jet}
\end{equation}
Here the endpoint terms at $N$ cancel exactly.  The moment integral tends
to $1$ with an explicit bound:
\begin{equation}
 \int_1^N b_j(t)\,dt
     =1-\frac1N\sum_{i=0}^j\frac{(\log N)^i}{i!},
 \qquad
 0\leq\frac1N\sum_{i=0}^j\frac{(\log N)^i}{i!}
       \leq\frac{2^{j+1}}{\sqrt N}.
 \label{eq:finite-log-moment}
\end{equation}
The first identity is proved by the finite recurrence
$U_0(N)=1-1/N$ and
$U_j(N)=U_{j-1}(N)-b_j(N)N$ for $j\geq1$.
For comparison, the corresponding recurrence for
$V_j(N)=\frac1{j!}\int_1^N t^{-3}(\log t)^jdt$ is
\[
 V_0(N)=\frac{1-N^{-2}}2,\qquad
 V_j(N)=\frac12V_{j-1}(N)
             -\frac{(\log N)^j}{2j!N^2}.
\]
This proves the coefficient estimate based on
$V_j(\infty)=2^{-j-1}$ by a bounded-index recurrence, without induction on
a formula quantifying over all real precisions.

By the definition of $u_j$, equations~\eqref{eq:finite-zeta-jet} and
\eqref{eq:finite-log-moment} give
\begin{equation}
 Z_j=1+u_j.
 \label{eq:zeta-jet-identification}
\end{equation}
For $j=0$ the contribution $\mathbf1_{j=0}$ is the constant $1$ already
included in $u_0$; for $j>0$ it is absent.

\paragraph{A finite rectangle gives the product coefficients.}
For a fixed $j$ the binomial identity and
$\log(uv)=\log u+\log v$ give the finite convolution identity
\begin{align}
 \sum_{i=0}^jZ_i(N)M_{j-i}(N)
  &=\sum_{u=1}^N\sum_{v=1}^N
       \frac{\chi(v)}{u^2v^2}
       \frac{(\log(uv))^j}{j!}.
 \label{eq:finite-moment-convolution}
\end{align}
The terms with $uv\leq N$ are exactly
\[
 c_j(N)=\sum_{n=1}^N\frac{a_n(\log n)^j}{j!n^2},
 \qquad a_n=\sum_{v\mid n}\chi(v).
\]
Every factor pair of $n\leq N$ occurs in the rectangle, so this last
assertion is an equality of bounded divisor sums.  The remaining terms have
$N<uv\leq N^2$, and their absolute sum is at most
\[
 \sum_{n>N}\frac{d(n)(\log n)^j}{j!n^2}
       \leq8\,4^jN^{-1/4}.
\]
The same bound controls $|c_j-c_j(N)|$.  Thus taking the explicitly bounded
tails in~\eqref{eq:finite-moment-convolution} proves
\begin{equation}
 c_j=\sum_{i=0}^j Z_i l_{j-i}
     =\sum_{i=0}^j u_i l_{j-i}+\sum_{i=0}^j l_i.
 \label{eq:finite-product-coefficient-identity}
\end{equation}
No rearrangement of an unbounded double series is needed: the exact
identity is on an $N$-by-$N$ rectangle, and all its discrepancies have
displayed positive tail bounds.

For clarity, one common bound is available.  To establish these identities
simultaneously for $j\leq K$ with final accuracy $2^{-r}$, first put
$p=r+2K+10$ and take
\[
 N=2^{4(p+K+8)}(q+1)^4 7^{4(K+1)}.
\]
Each of the unweighted moment tails, coefficient-integral tails, boundary
terms, and the bound in~\eqref{eq:finite-log-moment} is less than
$2^{-p-7}$.  The sum of the positive coefficient tails after weighting by
$(7/4)^j$ is also less than $2^{-p-7}$.  Passing from the truncated
moments to one convolution amplifies its individual tail errors by at most
$6(K+1)2^K\leq2^{2K+4}$.  These allowances are therefore smaller than
$2^{-r}$, with room for the additional rational sample errors.  Assign
those sample errors by the finite-operation budget before doing the
calculation.  Separate smaller cutoffs are possible, but unnecessary.

\paragraph{The endpoint coefficient sum.}
The remaining equality $\varepsilon=\sum_{i\geq0}l_i$ can also be checked
without a termwise interchange of unbounded operations.  Let
$E_H(x)=\sum_{i=0}^H x^i/i!$, with $E_{-1}=0$.  A finite sum of the integral
coefficient formulas gives
\[
 \sum_{i=0}^H l_i
 =\int_1^\infty A_\chi(t)t^{-3}
              \bigl(2E_H(\log t)-E_{H-1}(\log t)\bigr)dt.
\]
For $t\geq1$ both exponential polynomials are between $0$ and $t$.
Consequently the tail of this whole finite sum of integrals after $N$ is
at most $3q/N$, independently of $H$.  On $[1,N]$, choose $H$ large enough
that both $E_H$ and $E_{H-1}$ approximate $\exp$ to the assigned accuracy
on $[0,\log N]\subseteq[0,N]$.  The explicit exponential cutoff in the
foundations section supplies such an $H$, with $H$ at least any previously
specified coefficient index.  The integrand then approximates
$A_\chi(t)t^{-2}$, whose integral is $L(1,\chi)$ and whose tail is at most
$q/N$.  Both the interval and polynomial degree have thus been chosen by
primitive recursive bounds.

For a fixed $J$, split the finite sum at $J$ and use
$|l_i|\leq2q2^{-i}$ on its remaining coefficients.  The result at every
requested accuracy is
\[
 \left|\varepsilon-\sum_{i=0}^J l_i\right|
       \leq2q2^{-J}.
\]
It follows that $b_j=\sum_{i=0}^jl_i-\varepsilon$ has the bound
$|b_j|\leq2q2^{-j}$.  Substituting this definition into
\eqref{eq:finite-product-coefficient-identity} gives precisely
$c_j=\varepsilon+h_j$.

\paragraph{The induction invariant.}
Fix $q,K,r$ and the displayed cutoffs before beginning any of these loops.
The entries of the arrays are Gaussian rationals approximating
$Z_i(N),M_i(N),u_i,l_i$ and the partial integrals.  An invariant records
the stated finite sum or recurrence and its rational error allowance for
every index $i$ up to the current index.  This bounded conjunction is a
Boolean product.  The recurrence for the $V_i$ multiplies the previous
error by $1/2$ and adds the current local error; convolution has the finite
amplification bound just displayed.  The same finite-operation budget
handles powers of approximated logarithms.

Thus induction is used on equations and inequalities between coded
rationals, with bounded index assertions.  The accuracy $r$ is an external
number parameter in that induction, not an unbounded quantifier inside
its formula.  The resulting universally quantified error bounds are
obtained only after that quantifier-free induction is complete.  At the
final Dirichlet inequality the reserved positive slack lets one choose a
single sufficiently fine rational accuracy; there is no need to decide
equality of an arbitrary real name.
